\FloatBarrier


\section{Prefix Truncation}
\label{sec:truncation}

\subsection{The Hash Function}

ERE size grows with the bit-length of keys.
We need to compress keys into a smaller range $[0, 2^K)$ while preserving locality – so
that a query interval maps to a bounded number of intervals in the compressed space.
This need not be a hash in the cryptographic sense; the simplest locality-preserving
compression is plain bit-truncation:
\[
    h(x) = \lfloor x / 2^t \rfloor,
\]
keeping the top $K = L - t$ bits of each key (where $L$ is the key width in bits).

Properties of the hash:
\begin{itemize}
    \item \textbf{Locality-preserving:} truncation preserves order --- if $a < b$, then
    $h(a) \leq h(b)$.
    Intervals map to intervals.
    \item \textbf{Contiguous phantoms:} recall that the phantom points of a key $x$ are
    all non-keys $y$ with $h(y) = h(x)$ --- the pre-images that cause false positives.
    Under truncation they all fall within an interval of length~$2^t$ in the
    original universe: exactly the non-keys sharing the same $K$-bit prefix as $x$.
    This contrasts with the SODA hash, where phantoms are scattered uniformly.
\end{itemize}

\begin{figure}[h]
    \centering
    \includegraphics[width=0.85\textwidth]{phantom_comparison}
    \caption{Green strips show the \textbf{collision zone} of stored key $x$: starting
    positions $q$ of queries $[q, q+\mathcal{L}-1]$ that are false positives due to $x$
    (i.e.\ $x \notin [q,q+\mathcal{L}-1]$ yet $h(x)\in h([q,q+\mathcal{L}-1])$).
    Each scattered phantom of the random hash contributes $\mathcal{L}$ such positions
    (phantom at the strip's right edge); truncation's contiguous block limits the
    total to $2^t - 1$, independent of $\mathcal{L}$.}
    \label{fig:phantom-comparison}
\end{figure}

\begin{remark}
    Throughout this section, $2^t = 2^{L - K}$ is the universe compression factor; the number of
    original-universe points that hash to the same value in $[0, 2^K)$.
    For truncation this matches the definition: $t$ bits are dropped from the back of each key.
\end{remark}

\subsection{ARE Construction With Truncation}

Combining truncation with ERE yields a filter where the contiguous phantom structure
(Figure~\ref{fig:phantom-comparison}) changes the collision analysis.
We define the \textbf{collision zone} $\mathrm{CZ}(x)$ of stored key $x$ as the set of
query starting positions $q$ for which $[q, q+\mathcal{L}-1]$ is a false positive due
to $x$:
\[
    \mathrm{CZ}(x) \;=\; \bigl\{\, q \;:\; x \notin [q, q+\mathcal{L}-1] \text{ yet } h(x) \in h([q, q+\mathcal{L}-1]) \,\bigr\}.
\]
Its cardinality $|\mathrm{CZ}(x)|$ counts the false-positive starts contributed by $x$.

\begin{itemize}
    \item \textbf{SODA hash:} each of the $2^t - 1$ phantom positions is far from $x$,
    so all $\mathcal{L}$ starting positions that cover it also exclude $x$;
    $|\mathrm{CZ}(x)| = \mathcal{L}(2^t - 1)$ (Section~\ref{sec:soda}).
    \item \textbf{Truncation:} all $2^t - 1$ phantoms share one contiguous block with $x$.
    A query misses $x$ only when it starts inside $(x, \text{block\_right}]$ or ends
    inside $[\text{block\_left}, x)$; these two intervals have widths $(2^t - 1 - r)$
    and $r$ respectively, so $|\mathrm{CZ}(x)| = 2^t - 1$ regardless of $\mathcal{L}$
    (Eq.~\eqref{eq:trunc-fpr}).
\end{itemize}

The additive (rather than multiplicative) interaction eliminates the $\mathcal{L}$
factor from the filter's space formula:

\begin{table}[H]
    \centering
    \begin{tabular}{@{}llll@{}}
        \toprule
        Hash & $|\mathrm{CZ}(x)|$ & $K$ for FPR $\leq \varepsilon$ & Filter BPK \\
        \midrule
        Random (SODA) & $\mathcal{L} \times (2^t-1)$ & $\log_2(n\mathcal{L}/\varepsilon)$ &
        $\log_2(\mathcal{L}/\varepsilon) + O(1)$ \\
        Truncation & $2^t - 1$ & $\log_2(n/\varepsilon)$ &
        $\log_2(1/\varepsilon) + O(1)$ \\
        \bottomrule
    \end{tabular}
    \caption{Choosing truncation over the SODA hash reduces the resulting filter's BPK by
        $\log_2(\mathcal{L})$. At $\mathcal{L} = 128$, this is a saving of 7~bits per key.}
    \label{tab:trunc-vs-soda}
\end{table}

\begin{figure}[h]
    \centering
    \includegraphics[width=0.85\textwidth]{trunc_tradeoff_uniform}
    \caption{FPR vs BPK for truncation vs SODA on uniform keys ($n = 2^{20}$,
        $\mathcal{L} = 128$). The SODA curve is shifted right by
        ${\approx}\log_2(128) = 7$~BPK --- matching the theoretical prediction.}
    \label{fig:trunc-uniform}
\end{figure}

\FloatBarrier

\subsection{The Price of Contiguity}
\label{subsec:trunc-failure}

The same contiguity that saves space may also be a weakness in other scenarios.

\subsubsection{FPR for Uniform Key Spacing}
\label{subsubsec:trunc-fpr}

Let $R$ denote the uniform gap between consecutive keys.
Assume keys $x_0,\, x_0{+}R,\, x_0{+}2R,\ldots$
A query $[q,\, q{+}\mathcal{L}{-}1]$ is truly empty iff it fits inside a gap:
$q = x_i + s$, $s \in [1,\, R{-}\mathcal{L}]$.

A false positive requires some key's hash value to fall inside $h([q, q+\mathcal{L}-1])$.
Since the query is empty, every block \emph{fully covered} by the query is key-free:
any key inside such a block would lie in the query and contradict emptiness.
Only the two \emph{edge} blocks --- the one straddled by $q$ on the left and the one
straddled by $q + \mathcal{L} - 1$ on the right --- can hold a key whose hash collides
with $h([q, q+\mathcal{L}-1])$, and it suffices to consider the immediate neighbours
$x_i$ and $x_{i+1}$: order-preservation of $h$ gives $h(x_{i+2}) \geq h(x_{i+1})$, so
if $x_{i+2}$ collides with the query then $x_{i+1}$ does too --- further keys cannot
trigger \emph{new} false-positive events (symmetrically on the left).
The argument holds regardless of how many blocks the query spans.

Let $\alpha_i = x_i \bmod 2^t$ and $\alpha_{i+1} = x_{i+1} \bmod 2^t$ be the offsets of
the neighbours within their respective blocks.
A key's \textbf{phantom block} is the interval of points sharing the same $K$-bit prefix
as the key --- i.e., the $2^t$ universe positions that truncation maps to the same fingerprint.
\begin{itemize}
    \item \textbf{Left overlap:} the query overlaps $x_i$'s phantom block iff
    $q \leq x_i + 2^t - 1 - \alpha_i$, i.e., $s \leq 2^t - 1 - \alpha_i$ ---
    contributing $2^t - 1 - \alpha_i$ values of $s$.
    \item \textbf{Right overlap:} the query overlaps $x_{i+1}$'s phantom block iff
    $q + \mathcal{L} - 1 \geq x_{i+1} - \alpha_{i+1}$, i.e.,
    $s \geq R - \mathcal{L} + 1 - \alpha_{i+1}$ --- contributing $\alpha_{i+1}$ values of $s$.
\end{itemize}

When $R - \mathcal{L} \ge 2^{t+1}$, the two ranges of $s$ are disjoint, so the
false-positive count per gap is $(2^t - 1 - \alpha_i) + \alpha_{i+1}$.
Treating $\alpha_i, \alpha_{i+1}$ as uniform on $\{0, \ldots, 2^t - 1\}$ (mean $(2^t-1)/2$)
and dividing by the $R - \mathcal{L}$ empty-query positions:
\begin{equation}
    \mathrm{FPR}
    \;=\; \frac{2^t - 1}{R - \mathcal{L}}
    \;\approx\; \frac{2^t}{R - \mathcal{L}}.
    \label{eq:trunc-fpr}
\end{equation}

Capping at 1 and handling the $R \leq \mathcal{L}$ case:
\[
    \mathrm{FPR} =
    \begin{cases}
        0 & R \leq \mathcal{L} \text{ (no empty queries)}, \\
        \min\!\bigl(1,\ 2^t/(R - \mathcal{L})\bigr) & R > \mathcal{L}.
    \end{cases}
\]

Dense or sequential keys ($R \leq \mathcal{L}$) have FPR~$= 0$ for free: every query
of length $\mathcal{L}$ hits a key, so the filter is exact.
When $R > \mathcal{L}$, Eq.~\eqref{eq:trunc-fpr} reads $\mathrm{FPR} \approx 2^t / (R - \mathcal{L})$, so the gap excess $R - \mathcal{L}$ has to be a large multiple of the
phantom-block width $2^t$ to keep FPR low.
Concretely: at $R - \mathcal{L} = 10 \cdot 2^t$ the FPR is ${\approx}0.1$; at
$R - \mathcal{L} = 2 \cdot 2^t$ it is already ${\approx}0.5$; at the floor
$R - \mathcal{L} \leq 2^t$ the phantom block covers (almost) all of the gap and the
formula saturates at FPR~$=1$.

\subsubsection*{Breakdown thresholds for SOSD datasets}

From Eq.~\eqref{eq:trunc-fpr}, high FPR occurs when $2^t \gtrsim R - \mathcal{L}$,
so the critical truncation depth is $t^* = \lfloor\log_2(R - \mathcal{L})\rfloor$.
Table~\ref{tab:trunc-breakdown} shows $t^*$ for typical SOSD gaps and several query lengths.

\begin{table}[H]
    \centering
    \begin{tabular}{@{}llrrrr@{}}
        \toprule
        Dataset       & Typical $R$     & $\mathcal{L}=4$ & $\mathcal{L}=16$ & $\mathcal{L}=128$ & $\mathcal{L}=1024$ \\
        \midrule
        Wiki / Books  & 1--5            & $0^*$           & $0^*$            & $0^*$             & $0^*$              \\
        Facebook P50  & 28              & 4               & 3                & $0^*$             & $0^*$              \\
        Facebook mean & 387             & 8               & 8                & 8                 & $0^*$              \\
        OSM P50       & $4.5\times10^7$ & 25              & 25               & 25                & 25                 \\
        \bottomrule
        \multicolumn{6}{l}{\footnotesize ${}^{*}$~$R \leq \mathcal{L}$: every query contains a key.}
    \end{tabular}
    \caption{Critical truncation depth $t^*$ at which FPR~$\to 1$ for SOSD gap statistics
        (Table~\ref{tab:sosd-gaps}) and varying query length~$\mathcal{L}$.}
    \label{tab:trunc-breakdown}
\end{table}

Wiki and Books have gaps well below any practical $\mathcal{L}$, so every query hits a key.
Facebook's tail (gaps up to ${\sim}700$) breaks at $t \approx 8$.
OSM has $t^* \approx 25$ for all query lengths.
In a 64-bit universe this means keeping $K = 64 - 25 = 39$ bits just to avoid complete failure;
ERE compresses these to $K - \log n \approx 19$~BPK at $n = 2^{20}$.
Achieving FPR~$\leq 10^{-3}$ requires $t \leq 15$ ($K \geq 49$, $\approx 29$~BPK) under uniform queries.
Thus, truncation is impractical for clustered data.

\begin{remark}
    The derivation assumes \textbf{queries are placed uniformly at random in the universe},
    so $\alpha_i, \alpha_{i+1}$ are uniform on $\{0, \ldots, 2^t-1\}$.
    Benchmark queries (Section~\ref{sec:query-distribution}) are drawn from the gaps
    between keys, so they land near phantom block more often than a uniform draw would.
\end{remark}

The SODA hash avoids this entirely: $\mathrm{FPR} \leq \varepsilon$ for any key distribution, at the cost of
$\log_2(\mathcal{L})$ extra bits per key.

\begin{figure}[h]
    \centering
    \includegraphics[width=0.85\textwidth]{trunc_tradeoff_osm}
    \caption{FPR vs BPK on OSM ($n = 2^{20}$, $\mathcal{L} = 128$, in-gap queries, Section~\ref{sec:query-distribution}).
    Truncation is bounded below by a structural FPR floor set by the gap-to-$2^t$ ratio
        (Eq.~\ref{eq:trunc-fpr}): at 24~BPK it remains at FPR~$\approx 10^{-2}$
        while SODA reaches ${\sim}10^{-5}$.}
    \label{fig:trunc-osm}
\end{figure}

\FloatBarrier

\begin{remark}
Grafite~\citep{grafite2024} independently proposes an equivalent heuristic named
\emph{Bucketing} (their Section~4): keys are mapped via $x \mapsto \lfloor x/s \rfloor$
and the occupied bucket indices are encoded with Elias-Fano, with $s$ as the space-accuracy
trade-off parameter --- the same idea as prefix truncation with $s = 2^t$.
Our contribution relative to Grafite is the precise FPR formula
(Equation~\ref{eq:trunc-fpr}) and the per-dataset breakdown analysis
(Table~\ref{tab:trunc-breakdown}), which Grafite does not provide.
\end{remark}

\pagebreak


\section{Adaptive Filter}
\label{sec:adaptive}

The SODA hash compresses keys into a $K$-bit universe of size $2^K = n\mathcal{L}/\varepsilon$.
If the original key spread $\max(S) - \min(S)$ is already smaller than $2^K$,
no hash is needed at all --- we build ERE directly on the (normalized) keys
and get FPR~$= 0$ for free.
The adaptive filter checks this condition at build time.

\subsection{Exact Mode}
\label{subsec:exact-mode}

\begin{enumerate}
    \item \textbf{Normalize:} subtract $\min(S)$ from all keys, shifting the dataset to
    start at~0.
    \item \textbf{Measure spread:} $M = \lceil \log_2(\max(S) - \min(S)) \rceil$.
    \item \textbf{Decide:}
    \begin{itemize}
        \item $M \leq K$: the entire dataset fits in $K$ bits $\Rightarrow$ \textbf{exact
        mode}. Build ERE directly over normalized keys. No hash, FPR~$= 0$.
        \item $M > K$: data too spread $\Rightarrow$ \textbf{SODA mode}. Apply
        a hash drawn from a pairwise-independent family, same filter guarantees as Section~\ref{sec:soda}.
    \end{itemize}
\end{enumerate}

\subsection{When Does Exact Mode Trigger?}

Rewriting the condition
\[
    \max(S) - \min(S) < 2^{K} = \frac{n\mathcal{L}}{\varepsilon}
\]
in terms of
the average gap:
\[
    \bar{g} = \frac{\max(S) - \min(S)}{n - 1}
\]
\begin{equation}
    \bar{g} \;<\; \frac{\mathcal{L}}{\varepsilon}.
    \label{eq:exact-threshold}
\end{equation}
At $\mathcal{L} = 128$, $\varepsilon = 10^{-3}$ the threshold is $128{,}000$ ---
well above the Facebook median gap of~$28$ or the Books mean of~$5$,
so exact mode triggers on these datasets for free.

As we will see in Chapter~\ref{chap:hybrid}, this becomes especially powerful in combination with Seg-ARE:
each dense cluster is isolated into its own adaptive filter,
where the local spread is small enough for exact mode, giving FPR~$= 0$ on the dense
parts of the data.
